\textcolor{cyan}{In this work, we propose a} \st{We have used} crossed-parent cluster \textcolor{cyan}{protocol} to examine how composition and \st{cluster
geometry} \textcolor{cyan}{structural rearrangements} affect \ce{CO} adsorption in ten-atom Pt--Ge systems. In one route,
Ge atoms were introduced into a common \ce{Pt10} framework to generate
\ce{Pt6Ge4}, \ce{Pt5Ge5}, and \ce{Pt4Ge6} \textcolor{cyan}{that represent different relative metallic/oxophylic environments}. In the reverse route, the
\ce{Pt5Ge5} framework was converted into a pure \ce{Pt10} cluster.
{\textcolor{cyan}{Taken} together, these routes show how composition and geometry act jointly.
Replacing Pt by Ge changes the local chemical environment, bonding pattern,
and structural response of the cluster.}

AdNDP provides a simple bonding picture for these changes. Although the pure \(\ce{Pt10}\) derivative was generated from the reported global-minimum \(\ce{Pt5Ge5}\) framework, its optimized structure is described mainly by delocalized multicenter metallic bonding. In contrast, the Ge-containing derivatives favor more localized \(\ce{Pt-Ge}\) and \(\ce{Pt-Pt}\) bonding \textcolor{cyan}{schemes} \st{elements}. Within the representative \ce{Pt10 -> Pt_xGe_{10-x}}
crossed-parent sequence, the mean binding energies become progressively less
negative as the Ge content increases. The Ge-containing carbonyls also
generally show slightly longer Pt--C contacts and less activated \ce{C-O}
bonds.
These structural and vibrational changes are consistent with weaker
metal-to-\ce{CO} back-donation.

The affinity for \ce{CO} also depends strongly on the selected Pt site and on
cluster relaxation. Several initial adsorption positions converge to the same
minimum, whereas others give distinct carbonyls with different binding
energies. Calculations started from the previously reported global-minimum
structures further show that changing the parent isomer or spin state can
alter both the mean affinity and the range of binding energies. {The
initial ECoN descriptors show that Pt sites with a larger Ge contribution
generally bind \ce{CO} less strongly.} The exceptions identified
for site 9 of \ce{Pt6Ge4} and site 3 of the GM-parent \ce{Pt4Ge6} series show
that ECoN is best used as a description of the local environment rather than
as a direct predictor of the binding energy. Changes in ECoN upon adsorption
also describe local recoordination, but do not directly measure the energetic
cost of reorganizing the complete cluster.

\blue{The frozen-fragment analysis provides the corresponding energetic picture}.
For \ce{Pt6Ge4}, the more strongly bound carbonyl has the more stabilizing
frozen-fragment interaction, but its larger metal-reorganization penalty partly reduces
this advantage. For \ce{Pt4Ge6}, the frozen-fragment interaction and metal
reorganization both favor the more strongly bound carbonyl. In \ce{Pt5Ge5},
the favorable metal reorganization of the site-6/8 carbonyl outweighs its less
stabilizing frozen-fragment interaction and gives the more negative binding energy.
Its negative metal-reorganization term shows that the \ce{Pt5Ge5}
framework is stabilized at the geometry adopted in that carbonyl. For the
strongest crossed-parent \ce{Pt10CO} structure, the metal geometry
extracted from the complex relaxes to a lower-energy triplet \ce{Pt10}
minimum. This result shows that \ce{CO} adsorption at site 6 can promote
reorganization of \ce{Pt10} toward a lower-energy triplet isomer. Overall,
increasing the Ge content weakens the average \ce{CO} adsorption within the
sampled crossed-parent composition sequence\textcolor{cyan}{, in line with previous studies.\cite{ugartemendia:2022}}\st{, but} \textcolor{cyan}{Nonetheless, the present established protocol points to a more complex description of the CO adsorption mechanism, where} the magnitude and ordering of the \st{individual} binding \textcolor{cyan}{strength} \st{energies} \textcolor{cyan}{at indepedent active sites} reflect the combined effects of composition, bonding, parent geometry, spin state, and structural relaxation. \textcolor{cyan}{Namely, CO binding energies should not be interpreted solely as a local descriptor of the adsorption site, but rather a property of the full cluster--CO complex.}
